% ============================================================================
\section{Conclusiones}

Completado todo el análisis empírico, de simulación y de inferencia estadística planteado en el estudio, se formulan las conclusiones finales de acuerdo con el objetivo general y los objetivos específicos:

Respecto del objetivo general, la evaluación fuera de muestra sobre las revisiones del PBI peruano (\(n=250\), 2003--2024) demuestra que la computación de reservorio cuántico de cinco qubits supera puntualmente a la referencia nula y a la corrección clásica por media histórica (\(\text{RMSE}=0{,}478\) frente a \(0{,}513\), \(p=0{,}0175\) sin ajuste), pero \textbf{no mejora a los controles clásicos no lineales comparables} de igual dimensión de lectura: empata estadísticamente con la red clásica de estado eco ESN (\(\text{RMSE}=0{,}472\), \(p=0{,}730\)), es superada por características aleatorias de Fourier (\(\text{RMSE}=0{,}456\), con significancia tras ajustar por Holm--Bonferroni en el período previo a la pandemia, \(p_{\text{Holm}}=0{,}0265\)) y queda por debajo de los árboles potenciados acotados durante el choque pandémico (\(\text{RMSE}=0{,}349\) frente a \(0{,}363\)). En la familia confirmatoria global ningún contraste supera el umbral ajustado por multiplicidad (\(p_{\text{Holm}}\ge 0{,}1023\)), y el conjunto de confianza de modelos retiene conjuntamente a los competidores clásicos junto a QRC.

En relación con los objetivos específicos, se concluye que:
\begin{enumerate}
	\item \textbf{Propiedades estadísticas de las revisiones:} en la muestra depurada sin cambios de año base a horizonte continuo de once meses (\(n=350\)), las revisiones del PBI peruano presentan una estructura sistemática clara (Hipótesis~1), caracterizada por un sesgo positivo robusto de \(+0{,}347\) puntos porcentuales (\(p=5{,}3\times 10^{-8}\)), colas pesadas (curtosis en exceso de \(6{,}293\)), heterocedasticidad condicional, persistencia anual observable en tiempo real (\(\hat{\rho}_{12}=+0{,}386\)) y no linealidad concentrada únicamente en los años de pandemia 2020--2024.
	\item \textbf{Desempeño de referentes clásicos ante choques extremos:} antes de diciembre de 2019 (\(n=190\)), la regresión Ridge encabeza los métodos lineales tradicionales (\(\text{RMSE}=0{,}503\), \(R^2_{\text{OOS}}=+0{,}094\)); sin embargo, ante las caídas y rebotes atípicos de 2020--2021, las especificaciones lineales sin acotamiento extrapolan de forma excesiva (\(\text{RMSE}=0{,}731\) en el tramo final, quedando fuera del conjunto de confianza al \(75\)\%), mientras los métodos acotados (RFF, árboles potenciados y reservorios) evitan dicho deterioro.
	\item \textbf{Condiciones sintéticas y potencia estadística:} en doce escenarios calibrados sobre los parámetros observados del Perú, la región favorable para QRC resulta vacía (\(\mathcal{A}=\emptyset\)) y muestra mayor desventaja cuando la persistencia lineal es elevada (\(-20{,}90\)\% frente a Ridge). Asimismo, la muestra real (\(n=250\)) exige reducciones de error de al menos \(5{,}7\)\% a \(11{,}1\)\% frente a controles clásicos acotados (\(36{,}1\)\% frente a Ridge) para alcanzar una potencia del \(80\)\%, impidiendo confirmar ganancias pequeñas inferiores al \(2\)\%.
	\item \textbf{Valor económico condicionado:} al no cumplirse la condición previa de mejora predictiva robusta sobre la regresión Ridge y los reservorios clásicos equivalentes, no se ejecutó la simulación económica condicionada ni se atribuyó beneficio económico diferencial a la arquitectura cuántica.
\end{enumerate}

Cada conclusión queda respaldada por los cuadros comparativos, intervalos de confianza y pruebas de robustez presentados en este informe, evitando extrapolar resultados sintéticos a la economía peruana sin evidencia empírica adicional.
